% Bagian: computability
% Bab: computability-theory
% Subbagian: non-comp-set

\documentclass[../../../include/open-logic-section]{subfiles}

\begin{document}

\olfileid[id]{cmp}{thy}{ncp}
\olsection{Terdapat Himpunan yang Tidak Dapat Dihitung}


Kita telah melihat di atas bahwa setiap himpunan dapat dihitung terenumerasi
secara komputabel. Apakah kebalikannya benar? Hasil berikut menunjukkan bahwa
secara umum jawabannya tidak.

\begin{thm}\ollabel{thm:K-0}
Misalkan $K_0$ adalah himpunan
$\Setabs{\tuple{e, x}}{\cfind{e}(x) \fdefined}$. Maka $K_0$ terenumerasi
secara komputabel, tetapi tidak dapat dihitung.
\end{thm}

\begin{proof}
Untuk melihat bahwa $K_0$ terenumerasi secara komputabel, perhatikan bahwa
himpunan itu merupakan domain fungsi~$f$ yang didefinisikan oleh
\[
f(z) = \umin{y}{(\len{z} = 2 \land T((z)_0, (z)_1, y))}.
\]
Sebab, jika $\cfind{e}(x)$ terdefinisi, $f(\tuple{e, x})$ menemukan suatu
urutan komputasi yang berhenti; jika $\cfind{e}(x)$ tidak terdefinisi,
$f(\tuple{e, x})$ juga tidak terdefinisi; dan jika $z$ bahkan tidak mengodekan
suatu pasangan, $f(z)$ juga tidak terdefinisi.

Fakta bahwa $K_0$ tidak dapat dihitung tidak lain adalah ketakterputusan
masalah penghentian, \olref[hlt]{thm:halting-problem}.
\end{proof}

Himpunan $K_0$ adalah himpunan pasangan $\tuple{e,x}$ sedemikian sehingga
$\cfind{e}(x) \fdefined$, yakni $\tuple{e,x} \in K_0$ jika dan hanya jika
$\cfind{e}$ terdefinisi (berhenti) pada masukan~$x$. Karena itu, $K_0$ juga
disebut ``himpunan penghentian.'' Himpunan
$K = \Setabs{e}{\cfind{e}(e) \fdefined}$ adalah ``himpunan penghentian-diri.''
Himpunan ini sering digunakan sebagai himpunan kanonik yang tidak dapat
diputuskan.

\begin{thm}\ollabel{thm:K}
Himpunan penghentian-diri $K = \Setabs{e}{\cfind{e}(e) \fdefined}$ bersifat
!!{c.e.}, tetapi tidak dapat diputuskan.
\end{thm}

\begin{proof}
  Andaikan $K$ dapat diputuskan, yakni fungsi karakteristiknya $\Char{K}$
  dapat dihitung. Misalkan
  \[d(e) = \begin{cases}
  1 & \text{jika\/ $\Char{K}(e) = 0$}\\
  \fundefined & \text{selain itu.}
  \end{cases}
  \]
  Misalkan $k$ indeks~$d$, yakni $d \simeq \cfind{k}$. Maka
  $d(k) \simeq \cfind{k}(k)$. Hal ini bertentangan dengan fakta bahwa
  $d(k) \fdefined$ jika dan hanya jika $\cfind{k}(k) \fundefined$, yang
  mengikuti dari definisi~$d$.

  $K$ merupakan domain $f(x) = \umin{y}{T(x,x,y)}$, sehingga bersifat !!{c.e.}
\end{proof}

\end{document}

